% Contrato de prestación de servicios
% Compilar: /usr/bin/latexbuild contrato.tex
% Preámbulo compartido por el reporte, los manuales y el contrato.
% Se compila con XeLaTeX mediante /usr/bin/latexbuild.
\documentclass[11pt,a4paper]{article}

\usepackage{fontspec}
\setmonofont{JetBrainsMono Nerd Font Mono}[Scale=0.82]
\usepackage[spanish,es-nodecimaldot,es-noshorthands]{babel}

\usepackage[top=2.4cm,bottom=2.4cm,left=2.4cm,right=2.4cm]{geometry}
\usepackage{fancyhdr}
\usepackage{titlesec}
\usepackage{xcolor}
\usepackage{booktabs}
\usepackage{tabularx}
\usepackage{array}
\usepackage{enumitem}
\usepackage{float}
\usepackage{graphicx}
\usepackage{tikz}
\usetikzlibrary{positioning,arrows.meta}
\usepackage{listings}
\usepackage{xurl}

\makeatletter
\renewcommand*\l@section[2]{%
  \ifnum \c@tocdepth >\z@
    \addpenalty\@secpenalty
    \addvspace{1.0em \@plus\p@}%
    \setlength\@tempdima{2.6em}%
    \begingroup
      \parindent \z@ \rightskip \@pnumwidth
      \parfillskip -\@pnumwidth
      \leavevmode \bfseries
      \advance\leftskip\@tempdima
      \hskip -\leftskip
      #1\nobreak\hfil \nobreak\hb@xt@\@pnumwidth{\hss #2}\par
    \endgroup
  \fi}
\makeatother
\usepackage[hidelinks]{hyperref}

\definecolor{tnBlue}{HTML}{0EA5E9}
\definecolor{tnDark}{HTML}{0F172A}
\definecolor{codeGray}{HTML}{F4F6F8}
\definecolor{codeGreen}{HTML}{2E7D32}
\definecolor{codeRed}{HTML}{C62828}

\newcolumntype{Y}{>{\raggedright\arraybackslash}X}

\lstset{
  basicstyle=\ttfamily\footnotesize,
  breaklines=true,
  breakatwhitespace=false,
  frame=tb,
  framerule=0.4pt,
  rulecolor=\color{tnDark!40},
  framesep=5pt,
  keywordstyle=\color{tnBlue}\bfseries,
  commentstyle=\color{codeGreen}\itshape,
  stringstyle=\color{codeRed},
  showstringspaces=false,
  tabsize=2,
  columns=fullflexible,
  keepspaces=true,
  extendedchars=true,
  captionpos=t,
  xleftmargin=0pt,
  xrightmargin=0pt,
  aboveskip=0.9em,
  belowskip=0.9em,
  literate=
    {á}{{\'a}}1 {é}{{\'e}}1 {í}{{\'i}}1 {ó}{{\'o}}1 {ú}{{\'u}}1
    {Á}{{\'A}}1 {É}{{\'E}}1 {Í}{{\'I}}1 {Ó}{{\'O}}1 {Ú}{{\'U}}1
    {ñ}{{\~n}}1 {Ñ}{{\~N}}1 {ü}{{\"u}}1 {Ü}{{\"U}}1
    {¿}{{?`}}1 {¡}{{!`}}1
    {·}{{\textperiodcentered}}1
}

\pagestyle{fancy}
\fancyhf{}
\fancyhead[L]{\small Alfredo Ramírez Candia}
\fancyhead[R]{\small Diseño de Aplicaciones Web}
\fancyfoot[C]{\thepage}
\renewcommand{\headrulewidth}{0.4pt}

\titleformat{\section}
  {\Large\bfseries\color{tnDark}}{\thesection.}{0.5em}{}
\titleformat{name=\section,numberless}
  {\Large\bfseries\color{tnDark}}{}{0em}{}
\titleformat{\subsection}
  {\large\bfseries\color{tnBlue}}{\thesubsection}{0.5em}{}
\titlespacing*{\section}{0pt}{1.4em}{0.6em}
\titlespacing*{\subsection}{0pt}{1.05em}{0.35em}

\addto\captionsspanish{%
  \renewcommand{\tablename}{Tabla}%
  \renewcommand{\figurename}{Figura}%
  \renewcommand{\lstlistingname}{Listado}%
}

\newcommand{\datosPdf}[1]{%
  \hypersetup{%
    pdftitle={#1},%
    pdfauthor={Alfredo Ramírez Candia},%
    pdflang={es}%
  }%
}

\newcommand{\portada}[3]{%
  \begin{titlepage}
    \centering
    \vspace*{0.8cm}
    {\large\bfseries BENEMÉRITA UNIVERSIDAD AUTÓNOMA DE PUEBLA\par}
    \vspace{0.28cm}
    {\large FACULTAD DE CIENCIAS DE LA COMPUTACIÓN\par}
    \vspace{0.28cm}
    {\large LICENCIATURA/INGENIERÍA EN CIENCIAS DE LA COMPUTACIÓN\par}
    \vspace{1.35cm}
    \rule{\linewidth}{0.5pt}\par
    \vspace{0.45cm}
    {\Huge\bfseries #1\par}
    \vspace{0.38cm}
    {\Large #2\par}
    \vspace{0.32cm}
    {\large Asignatura: Diseño de Aplicaciones Web\par}
    \vspace{0.45cm}
    \rule{\linewidth}{0.5pt}\par
    \vspace{1.5cm}
    \begin{tabular}{rl}
      \textbf{Docente:}   & Miguel Ángel Vargas Lomelí \\[0.55cm]
      \textbf{Alumno:}    & Alfredo Ramírez Candia \\[0.55cm]
      \textbf{Matrícula:} & 202375947 \\[0.55cm]
      \textbf{Fecha:}     & 22 de noviembre de 2026 \\
    \end{tabular}
    \vfill
    {\small #3\par}
  \end{titlepage}
  \setcounter{page}{2}
}

\setcounter{tocdepth}{2}
\setcounter{secnumdepth}{2}
\setlist[enumerate]{leftmargin=1.5em,itemsep=0.45em,topsep=0.4em}
\setlist[itemize]{leftmargin=1.4em,itemsep=0.28em,topsep=0.35em}
\emergencystretch=2em
\urlstyle{tt}
\renewcommand{\arraystretch}{1.18}

\datosPdf{Contrato de prestación de servicios. DevStudio}

\begin{document}
\portada{Contrato de servicios}{Desarrollo de una aplicación web del tamaño de DevStudio}{Modelo académico, no constituye asesoría legal}

\tableofcontents
\clearpage

\section{Advertencia}

Este contrato es un modelo para la materia Diseño de Aplicaciones Web. Fija el
alcance y el costo de construir un software de la medida de DevStudio. No
sustituye la revisión de un profesional del derecho ni obliga a un cliente
real.

\section{Partes}

El contrato número PF-2026-001 se celebra en Puebla, México, el 22 de
noviembre de 2026.

\textbf{Prestador.} Alfredo Ramírez Candia, alumno de la Facultad de Ciencias
de la Computación de la Benemérita Universidad Autónoma de Puebla, matrícula
202375947. Domicilio para efectos de este modelo: Puebla, México. Correo de
contacto del proyecto: \texttt{contacto@devstudio.local}.

\textbf{Cliente.} DevStudio, negocio de desarrollo para comercios de Puebla,
en cuyo nombre se encarga el sistema descrito en este documento. Domicilio:
Puebla, México.

Las partes declaran tener capacidad para firmar este modelo con fines
académicos.

\section{Objeto}

El prestador diseña, construye y entrega una aplicación web de la medida de
DevStudio, compuesta por sitio público, portal de clientes, panel de
administración, API PHP y base SQL Server, más la documentación de usuario,
la documentación técnica y este contrato.

\section{Alcances}

\subsection{Módulos incluidos}

\begin{table}[H]
  \centering
  \caption{Lo que el precio de la sección 8 cubre.}
  \label{tab:alcance}
  \small
  \begin{tabularx}{\textwidth}{@{}l Y@{}}
    \toprule
    Módulo & Queda incluido \\
    \midrule
    Sitio público & Inicio, soluciones, casos, forma de trabajo, contacto y diagnóstico \\
    Acceso & Registro, inicio y cierre de sesión, roles cliente y administrador \\
    Portal & Perfil, foto, proyectos, avance, solicitudes y calendario \\
    Administración & Usuarios, clientes, servicios, proyectos, tecnologías, solicitudes y mensajes \\
    API de seguimiento & Actualizar el avance y confirmar o rechazar reuniones \\
    Datos & Esquema, semilla y migraciones idempotentes \\
    Entorno & Docker con Nginx, PHP 8.3 y SQL Server 2022 \\
    Documentación & Manual de usuario, manual técnico y este contrato, en PDF \\
    \bottomrule
  \end{tabularx}
\end{table}

\subsection{Entregables}

\begin{enumerate}
  \item Código fuente en repositorio y en archivo ZIP, sin secretos.
  \item \path{schema.sql}, \path{seed.sql} y las migraciones.
  \item Cuatro PDF generados con LaTeX.
  \item Sesión de demostración con las cuentas locales del seed.
\end{enumerate}

\subsection{Aceptación técnica}

El sitio se usa en móvil, tableta y escritorio. El cliente solo ve su
información. El administrador gestiona el contenido. La base se reconstruye
desde los scripts. Un cliente no puede escribir el avance ni responder una
reunión. La documentación permite instalar el sistema y recorrerlo.

\section{Limitaciones}

No entran en el precio, salvo un anexo firmado:

\begin{itemize}
  \item Pagos en línea, facturación o nómina.
  \item Aplicaciones nativas para iOS o Android.
  \item Chat, videollamada o notificaciones push.
  \item Integración con un CRM, un ERP u otro sistema del cliente.
  \item Correo de recuperación de contraseña o correos transaccionales.
  \item Pantalla nueva para que el desarrollador confirme reuniones: la acción existe en la API.
  \item Hosting, dominio, certificados y costos de nube.
  \item Soporte posterior a la garantía, auditoría de intrusión o cumplimiento normativo.
\end{itemize}

\section{Obligaciones}

El cliente entrega textos, datos de contacto y un responsable de las
revisiones. Responde en un plazo de cinco días hábiles y prueba el sistema en
el plazo de aceptación.

El prestador construye lo descrito, comunica el avance, entrega código y
documentación, hace la demostración y corrige durante la garantía los defectos
que sean del desarrollo.

\section{Cronograma}

\begin{table}[H]
  \centering
  \caption{Fases de un encargo de este tamaño.}
  \begin{tabularx}{\textwidth}{@{}l Y l@{}}
    \toprule
    Fase & Resultado & Cierre de referencia \\
    \midrule
    1 & Requisitos, mapa, guiones y navegación & 13 de noviembre de 2026 \\
    2 & Interfaz y modelo de datos & 15 de noviembre de 2026 \\
    3 & Portal, calendario, avance y foto & 22 de noviembre de 2026 \\
    4 & Pruebas en 375, 768 y 1280 píxeles & 22 de noviembre de 2026 \\
    5 & PDF y ZIP & 22 de noviembre de 2026 \\
    \bottomrule
  \end{tabularx}
\end{table}

Si el cliente no entrega insumos, los plazos se recorren por escrito.

\section{Precio}

El costo de hacer un software de esta medida es de \textbf{\$235{,}000 MXN}
más IVA. El desglose corresponde a los módulos de la sección 3 y no a un plan
comercial que el sitio ofrezca a un negocio.

\begin{table}[H]
  \centering
  \caption{Desglose del costo.}
  \label{tab:costo}
  \begin{tabular}{lr}
    \toprule
    Concepto & Importe (MXN) \\
    \midrule
    Modelación y diseño & 25{,}000 \\
    Sitio público, mobile first & 45{,}000 \\
    Autenticación y roles & 20{,}000 \\
    Portal, avance, calendario y perfil & 55{,}000 \\
    Panel de administración y API & 50{,}000 \\
    Base de datos y entorno Docker & 25{,}000 \\
    Documentación y transferencia & 15{,}000 \\
    \midrule
    Total, antes de IVA & 235{,}000 \\
    \bottomrule
  \end{tabular}
\end{table}

El pago se divide en tres partes: 40\,\% a la firma (\$94{,}000), 30\,\% al
aprobar la versión funcional (\$70{,}500) y 30\,\% contra entrega y aceptación
(\$70{,}500). La forma de pago de un encargo real sería transferencia. Este
modelo no incluye datos bancarios ni factura. El precio no cubre hosting ni
servicios de terceros.

\section{Cambios de alcance}

Cualquier función que no esté en la sección 3 se cotiza aparte. El trabajo
empieza cuando el cliente acepta esa cotización por escrito.

\section{Propiedad y confidencialidad}

Liquidado el total, el código y la documentación específicos del proyecto pasan
al cliente. El prestador conserva componentes genéricos y el derecho a mostrar
el trabajo en su portafolio, salvo que el cliente pida reserva por escrito.
La confidencialidad dura dos años después del fin del contrato, excepto la
información que ya sea pública o que pida una autoridad.

\section{Garantía y aceptación}

Durante 30 días naturales después de la aceptación, el prestador corrige sin
costo los defectos reproducibles que vengan del desarrollo. No cubre cambios
nuevos, modificaciones de terceros, fallas de red o de hosting, ni textos
aportados por el cliente.

El cliente tiene 10 días hábiles para aceptar o enlistar defectos. Sin
respuesta en ese plazo, el modelo da el entregable por aceptado.

\section{Terminación y ley aplicable}

Cualquiera puede terminar el contrato con un aviso de 15 días. Se paga el
trabajo entregado hasta esa fecha. El contrato se rige por las leyes de los
Estados Unidos Mexicanos. Para controversias, las partes acuden a los
tribunales de Puebla y renuncian a otro fuero.

\section{Firmas}

\vspace{1.1cm}
\noindent
\begin{minipage}[t]{0.47\textwidth}
  \textbf{Prestador}\\[1.5cm]
  Alfredo Ramírez Candia\\[0.8cm]
  Firma: \rule{3cm}{0.4pt}\\[0.6cm]
  22 de noviembre de 2026
\end{minipage}\hfill
\begin{minipage}[t]{0.47\textwidth}
  \textbf{Cliente}\\[1.5cm]
  DevStudio\\[0.8cm]
  Firma: \rule{3cm}{0.4pt}\\[0.6cm]
  22 de noviembre de 2026
\end{minipage}

\vspace{1.3cm}
\noindent\textbf{Testigos}

\vspace{0.8cm}
\noindent
\begin{minipage}[t]{0.47\textwidth}
  Nombre: \rule{3cm}{0.4pt}\\[0.8cm]
  Firma: \rule{3cm}{0.4pt}
\end{minipage}\hfill
\begin{minipage}[t]{0.47\textwidth}
  Nombre: \rule{3cm}{0.4pt}\\[0.8cm]
  Firma: \rule{3cm}{0.4pt}
\end{minipage}

\end{document}
